Paper~0 is the geometric foundation of this series.  It constructs the space in
which the elementary arithmetic operations act as motions, the dual reading of
those motions, the commutativization against which their order defect is
measured, and the contact structure of that defect.  It fixes explicit left- and
right-expanded conventions, proves the distributive expansion that relates two
expansions sharing an operator, and treats reciprocals, poles, the projective
point at infinity, affine and reciprocal fixed points, and the common
$2\times2$ matrix representation.

The present paper begins one level earlier and stays there.  Its first problem
is intrinsic: which finite binary expression trees have a unique legal order of
internal evaluation, including mixed operand-slot patterns?  The answer is the
single-spine classification.  Choosing an accumulator then gives the marked
spinal histories used throughout this paper.  We retain that terminology and
formalism unchanged.

The division of labour is therefore the following.
\begin{itemize}
  \item \textbf{Paper 0:} expansions, arithmetic motions, the models
  $\mathfrak E_0$ and $\mathfrak E_1$, ripple geometry, affine cocycles, the
  accumulative commutative space and tearing, and the contact structure of the
  tearing.
  \item \textbf{Paper I:} intrinsic sequential and branching histories, bilateral
  projective evaluation, the Scott--Lawson order topology of histories, the
  analysis three pillars (differential, integral, series), and the hole--cap
  fine geometry of the regular and singular zero theory.
\end{itemize}
Per the amendment recorded in
\path{governance/00d-paper-I-calculus-and-capping-amendment.md}, the analysis
and the hole--cap geometry are now developed \emph{here}; Paper~0's contact,
ACS, and punctured-AES chapters are the canonical upstream, and this paper
restates the specific primitives it needs at the same labels.

This paper imports, with labels and proofs unchanged from Paper~0:
\begin{enumerate}
  \item the regular arithmetic expression space
  $\bigl(M,g,a;\mu,\lambda\bigr)$ and its eikonal axiom
  $|\nabla a|_g^2=\mu^2+\lambda^2 a^2$ with $\mu\ne0$
  (\texttt{def:regular-aes}), together with the canonical arithmetic frame
  (\texttt{prop:canonical-arithmetic-frame});
  \item the arithmetic contact model $\mathcal C=\mathbb R^3_{(u,v,a)}$, the
  contact one-form $\alpha=da-\mu\,du-\lambda a\,dv$
  (\texttt{eq:arithmetic-contact-form}), the horizontal fields
  $D_u=\partial_u+\mu\partial_a$, $D_v=\partial_v+\lambda a\partial_a$
  (\texttt{eq:horizontal-arithmetic-fields}), and the contact curvature bracket
  $[D_u,D_v]=\mu\lambda\,\partial_a$ (\texttt{thm:contact-curvature});
  \item the accumulative commutative space $\mathcal A=\mathbb R^2_{(A,M)}$ and
  the weighted torsion--Stokes theorem (\texttt{thm:torsion-stokes});
  \item the punctured AES (\texttt{def:p0-punctured-aes}) and the two-sided
  collar contract (\texttt{def:p0-collar-contract}).
\end{enumerate}
These imports are restated in the chapters that use them; no theorem of Paper~0
is re-proved here, and the ordinary-admissibility / projective-continuation
distinction remains active throughout.
